% Part: lambda-calculus
% Chapter: introduction
% Section: fixed-point-combinator

\documentclass[../../../include/open-logic-section]{subfiles}

\begin{document}

\olfileid[hi]{lam}{int}{fix}
\olsection{अचल-बिंदु संयोजक}

मान लें आपके पास कोई लैम्ब्डा पद $g$ है और आप ऐसा दूसरा पद $k$ चाहते हैं
जिसमें $k$, ~$gk$ के $\beta$-तुल्य हो। हमारे संकेत-लिपि नियमों का उपयोग
करते हुए पद
\[
\fn{diag}(x) = xx
\]
और
\[
l(x) = g(\fn{diag}(x))
\]
परिभाषित करें; दूसरे शब्दों में, $l$ पद $\lambd[x][g(xx)]$ है। $k$ पद
$ll$ हो। तब
\begin{align*}
k & = (\lambd[x][g(xx)])(\lambd[x][g(xx)]) \\
& \red  g((\lambd[x][g(xx)])(\lambd[x][g(xx)])) \\
& = gk.
\end{align*}
यदि हम
\[
Y = \lambd[g][((\lambd[x][g(xx)])(\lambd[x][g(xx)]))]
\]
लें, तो $Yg$ और $g(Yg)$ अपचयित होकर किसी समान पद में पहुँचते हैं; अतः $Yg
\equiv_\beta g(Yg)$। इसे ``करी का संयोजक'' कहते हैं। यदि इसके बदले हम
\[
Y = (\lambd[xg][g(xxg)])(\lambd[xg][g(xxg)])
\]
लें, तो वास्तव में $Yg$ अपचयित होकर $g(Yg)$ बनता है, जो अधिक प्रबल कथन
है। $Y$ के इस दूसरे रूप को ``ट्यूरिंग का संयोजक'' कहते हैं।

\end{document}
